\cleardoublepage
\chapter{PENUTUP}
\label{chap:penutup}

\section{Kesimpulan}

Berdasarkan hasil perancangan, implementasi, dan evaluasi pada Bab IV
hingga Bab VI, tugas akhir ini menyimpulkan hal-hal berikut, disusun
sesuai urutan rumusan masalah dan tujuan pada Bab I:

\begin{enumerate}
	\item Model drain baterai non-linear bus listrik TransJogja berhasil
	dibangun berdasarkan Persamaan 1 \textcite{mamarikas2025}, dan
	divalidasi secara kualitatif terhadap baseline Rute L-1. Validasi
	menunjukkan kebutuhan pengisian daya bergeser dari 2 sesi per hari
	(kondisi riil) menjadi 3 sesi per hari (hasil simulasi), sebuah
	deviasi yang diterima secara sadar sebagai konsekuensi mempertahankan
	integritas model literatur, sebagaimana dijelaskan pada Bab IV dan
	Bab V.
	\item Mesin simulasi operasi armada berbasis peristiwa diskrit berhasil
	dibangun dan mampu menjalankan seluruh 20 koridor TransJogja secara
	simultan, termasuk mekanisme antrean di SPKLU dengan kapasitas
	terbatas, sebagaimana dijelaskan pada Bab V.
	\item Perbandingan skenario lokasi SPKLU (S15: 1 lokasi, S16: 2 lokasi,
	S17b: 4 lokasi, dengan total kapasitas gun yang setara) menunjukkan
	bahwa menambah jumlah lokasi saja tidak menjamin perbaikan headway --
	S16 justru lebih buruk dari S15 akibat pemecahan kapasitas gun di
	lokasi yang kurang tepat. Cara kapasitas didistribusikan antar lokasi
	ternyata lebih menentukan daripada jumlah lokasinya sendiri: pada
	jumlah lokasi dan total kapasitas yang sama, distribusi proporsional
	terhadap beban rute riil (skenario S17b) memberikan hasil yang jauh
	lebih baik daripada distribusi yang merata tanpa mempertimbangkan
	beban (skenario S17).
	\item Perbandingan skenario kapasitas SPKLU menunjukkan titik patah
	(\textit{regime shift}) yang tajam pada rentang kapasitas tertentu
	(Bab VI), bukan penurunan performa yang bertahap. Skenario dengan
	penambahan kapasitas 12\% di atas alokasi proporsional dasar (S28)
	memberikan tingkat degradasi rute mendekati nol.
	\item Perbandingan skenario mekanisme pengisian daya menunjukkan bahwa
	strategi penambahan infrastruktur fisik dan strategi algoritma
	penjadwalan cerdas menjawab persoalan yang berbeda: infrastruktur
	memperbaiki kepatuhan terhadap target headway, sedangkan algoritma
	memperbaiki keselamatan armada dengan biaya lebih rendah.
	Untuk evaluasi kepatuhan headway operasional harian, \textbf{S28}
	(26 gun + konfigurasi proporsional, tanpa algoritma dinamis) adalah
	skenario dengan performa terbaik yang telah diuji secara langsung.
	Untuk ketahanan operasional jangka panjang multi-hari, \textbf{S33}
	(infrastruktur S28 + algoritma V1) adalah satu-satunya skenario yang
	terbukti tangguh selama 14 hari beruntun dengan nol rute terdegradasi,
	meskipun dengan sedikit trade-off headway harian. Eksplorasi titik
	investasi minimum (skenario S35, 23 gun + algoritma V1) menemukan
	\textit{sweet spot} yang menghasilkan degradasi nol dengan kapasitas
	infrastruktur lebih rendah, namun skenario ini bersifat verifikasi
	tambahan di luar cakupan inti tugas akhir (Bab I, Batasan Masalah)
	dan memerlukan pengujian ketahanan multi-hari tersendiri sebelum dapat
	direkomendasikan secara operasional. Terkait kebutuhan penambahan
	armada, temuan skenario S26 menunjukkan bahwa menambah bus tanpa
	menambah kapasitas SPKLU secara setara justru kontraproduktif --
	degradasi rute meningkat dari 3,35 menjadi 10,60 rute per hari --
	sehingga penambahan armada baru direkomendasikan setelah kapasitas
	infrastruktur memadai, bukan sebagai solusi mandiri.
	\item Dasbor visual berhasil dibangun dan mampu menampilkan peta
	skematik rute serta animasi pergerakan bus per skenario, sebagai alat
	bantu komunikasi hasil kepada pemangku kepentingan, sebagaimana
	dijelaskan pada Bab V.
\end{enumerate}

Secara keseluruhan, dari seluruh skenario yang diuji, dengan kriteria ketat
(deviasi 0\%) tidak ada satu skenario pun yang meloloskan mayoritas dari 20
koridor terhadap target headway Dishub DIY; skenario terbaik hanya meloloskan
4 dari 20 rute (20\%). Namun, dengan toleransi operasional yang lebih realistis
($\pm$10 menit dari target), S28 meloloskan 17 dari 20 rute, menjadikannya
Pareto-dominan dibandingkan seluruh skenario alternatif yang diuji. Temuan ini
menunjukkan bahwa elektrifikasi penuh TransJogja dengan target headway yang
berlaku saat ini memerlukan kombinasi investasi infrastruktur dan strategi
operasional yang cermat, dan membuka pertanyaan lebih lanjut mengenai realisme
target headway tersebut untuk armada yang sepenuhnya elektrik (Bab VI).


\section{Saran}

Berdasarkan batasan masalah (Bab I) dan temuan pada Bab VI, beberapa
saran diajukan untuk pengembangan lebih lanjut:

\begin{enumerate}
	\item \textbf{Validasi lapangan pada skala lebih besar.} Validasi model
	drain baterai pada tugas akhir ini hanya menggunakan baseline Rute L-1
	(2 bus). Pengumpulan data operasional dari lebih banyak rute, seiring
	perluasan elektrifikasi TransJogja berjalan, akan memungkinkan
	validasi model yang lebih kuat secara statistik.
	\item \textbf{Optimasi lokasi SPKLU berbasis algoritma.} Tugas akhir ini
	membandingkan skenario lokasi yang ditentukan sebagai \textit{input}
	(Bab I, Batasan Masalah). Penelitian lanjutan dapat mengeksplorasi
	algoritma optimasi lokasi otomatis (misalnya metaheuristik atau
	\textit{reinforcement learning}) untuk mencari konfigurasi lokasi dan
	kapasitas yang lebih optimal secara sistematis, dibandingkan
	pendekatan perbandingan skenario yang digunakan pada tugas akhir ini.
	\item \textbf{Optimasi gabungan infrastruktur dan algoritma.} Skenario
	S35 pada tugas akhir ini menguji algoritma dinamis di atas
	infrastruktur proporsional yang sudah ditentukan secara sekuensial (23 gun), bukan
	mengoptimasi keduanya secara bersamaan sejak awal. Penelitian lanjutan
	dapat mengeksplorasi apakah konfigurasi infrastruktur yang berbeda
	memberikan hasil yang lebih baik ketika diuji bersama
	algoritma pengisian daya dinamis sejak tahap perancangan.
	\item \textbf{Perbaikan algoritma \textit{queue-aware}.} Algoritma V2
	pada tugas akhir ini menunjukkan fenomena osilasi sistemik akibat
	respons serentak antar-bus terhadap informasi antrean bersama (Bab
	VI). Penelitian lanjutan dapat mengeksplorasi mekanisme yang mencegah
	respons serentak ini, misalnya melalui pengacakan waktu respons atau
	koordinasi terpusat antar-bus.
	\item \textbf{Simulasi multi-hari dengan status baterai berkelanjutan.}
	Tugas akhir ini menggunakan simulasi harian yang berdiri sendiri
	(\textit{single-day/stateless}), dengan satu pengujian tambahan 14
	hari sebagai verifikasi ketahanan di luar cakupan inti (Bab VI).
	Penelitian lanjutan dapat mengembangkan simulasi multi-hari penuh
	sebagai bagian dari cakupan inti, termasuk pemodelan kebijakan
	pengisian daya di \textit{pool} pada malam hari secara lebih rinci.
	\item \textbf{Pemodelan variasi beban penumpang dan kondisi cuaca.}
	Tugas akhir ini tidak memodelkan pengaruh beban penumpang maupun
	cuaca terhadap konsumsi energi (Bab I, Batasan Masalah). Penelitian
	lanjutan dapat mengintegrasikan variabel-variabel ini ke dalam model
	drain baterai untuk meningkatkan akurasi estimasi konsumsi energi.
	\item \textbf{Diskusi realisme target headway dengan Dishub DIY.}
	Mengingat tidak ada skenario yang meloloskan mayoritas rute terhadap
	target headway yang berlaku saat ini (Bab VI), disarankan agar hasil
	tugas akhir ini didiskusikan dengan Dishub DIY untuk mengevaluasi
	apakah target tersebut perlu disesuaikan khusus untuk armada yang
	sepenuhnya elektrik, atau apakah investasi infrastruktur dan strategi
	operasional perlu ditingkatkan lebih lanjut dari rekomendasi pada
	tugas akhir ini.
\end{enumerate}